\lecture{13}{Wed 05 Oct 2022 10:37}{}

\begin{remark}[Commutator]
	$[x,y] = x y x^{-1} y^{-1} $ gives a measure of how nonabelian a group is.
\end{remark}


\begin{theorem}(6.1.9)
	$A_n$ is simple for $n\ge 5$.
\end{theorem}
\begin{proof}
	The last theorem proves $A_5$ is simple. We prove by induction, and assume that $A_m$ is simple for $5\le m<n$, and seek to prove that $A_n$ is simple.

	Suppose $\{e\} \neq N \triangleleft A_n$, and try to show that $N$ contains an element $\sigma$ which leaves an element of $\{1,2,\ldots,n\} $  fixed. Suppose that no such $\sigma$ exists. Then there exists distinct elements $a,b,c,d \in \{1,2,\ldots,n\} $ such that \[
		\sigma(a) = b, \sigma(c) = d.
	.\] 
	Since $n\ge 6$, there is an element $\tau \in A_n$ with \[
		\tau = (a,b)(c,d,e, f)
	.\] 
	for $e,f \not\in \{a,b,c,d\} $ with $e \neq f$.
	By normality of $N$ in $A_n$, one has $\tau^{-1} \sigma\tau \in N$.
	Then \begin{align*}
		\tau^{-1} \sigma\tau(b) &= \tau^{-1} \sigma(a) = \tau^{-1} (b) = a\\
		\tau^{-1} \sigma\tau(d) &= \tau^{-1} \sigma(e) = \tau^{-1} (e) = d
	.\end{align*} 
	Then
	\begin{align*}
		\tau^{-1} \sigma\tau\sigma(a) &= \tau^{-1} \sigma\tau(b) = a\\
		\tau^{-1} \sigma\tau\sigma(c) &= \tau^{-1} \sigma\tau(d) = d\\
	\end{align*}
	So $\tau^{-1} \sigma\tau\sigma$ fixes $a$ and moves $c$, so it is not the identity, and there is an element of $N$ fixing  $a$. We have a contradiction.

	We can suppose, by relabeling, that $\sigma \in N\setminus \{\text{id}\} $ and $\sigma(1) = 1$. The elements of $A_n$ which fix $1$ form a subgroup of $A_n$, say $H_1 \le A_n$, with $H_1 \cong A_{n-1}$. 
	Then $N \cap H_1 \neq \{\text{id}\} $ (since $\sigma \in N \cap H_1$), and $N \cap H_1 \triangleleft H_1$.
	But $H_1$ is simple (because $A_{n-1}$ is simple), so $N \cap H_1 = H_1$. 
	Then $H_1\le N$.

	Let $\tau$ be any permutation with $\tau(2) = 1$. Then $\tau^{-1} \sigma\tau(2) = \tau^{-1} \sigma(1) = \tau^{-1} (1) = 2$.
	So $H_2 := \tau^{-1} H_1\tau$ is a subgroup of $A_n$ that fixes $2$.
	By normality of $N$ in $A_n $ we have $\tau^{-1} H_1 \tau \le N$, and $\tau^{-1} H_1\tau \cong A_{n-1 }$.
	Also, \[
		H_1 \cap \tau H_1 \tau \cong A_{n-2}
	\],
	since it fixes both $1$ and $2$.

	Note, since $H_1 \cap H_2$ fixes both $1$ and $2$, so $H_1 \cap H_2 \cong A_{n-2}$.
	Consider the group generated by $H_1$ and $H_2$ :
	\[
		\text{gp}(H_1,H_2) 
		=\{g_1,\ldots,g_k: g_i \in H_1 \lor g_i \in H_2\} .
	\]
	Observe that $\text{gp}(H_1, H_2)\le \mathbb{N}$ since $H_1 \le N$ and $H_2 \le N$. 
	Also, the set $H_1H_2$ is contained in $\text{gp}(H_1, H_2)$ and hence in $N$.

	We compute $|H_1H_2|$. We look at how many representations an element $n \in H_1H_2$ has in the shape $n = h_1h_2$, with $h_1 \in H_1$ and $h_2 \in H_2$.

	If $n = g_1g_2$ and $n=h_1h_2$, $g_1,h_1 \in H_1$, $g_2, h_2 \in H_2$, then $g_1g_2 = h_1h_2$.
	So $H_1 \ni h_1^{-1} g_1 = h_2 g_2^{-1} \in H_2$, and hence, for some $t \in H_1 \cap H_2$, we have $h_1^{-1} g_1 = t = h_2 g_2^{-1}$.

	Then the number of ways that $n \in H_1H_2$ is represented as a product $h_1h_2$ with $h_1 \in H_1$ and $h_2 \in H_2$ is at most $|H_1 \cap H_2|$. Then \[
		|H_1H_2| \ge |H_1| |H_2| / |H_1 \cap H_2|
	.\] 
	Hence 
	\begin{align*}
		|N| &\ge  |H_1 H_2| \ge |H_1| |H_2| / |H_1 \cap  H_2| \\
		    &= \frac{|A_{n-1}||A_{n-1}|}{|A_{n-2}|} \\
		    &= \frac{\frac{1}{2}(n-1)!\frac{1}{2}(n-1)!}{\frac{1}{2}(n-2)!}\\
		    &= \frac{1}{2} (n-1) (n-1)!
	.\end{align*} 
But $N \triangleleft A_n$, so $|N| \le \frac{1}{2} n!$. Thus $1 \le |A_n / N| \le \frac{\frac{1}{2} n!}{\frac{1}{2} (n-1) (n-1)!} = \frac{n}{n-1}<2$ when $n >2$. But now $n \ge 6$, so $|A_n / N| = 1$ and $N = A_n$. Thus $A_n$ is simple for $n \ge 5$.
\end{proof}

